\documentclass[11pt]{article}
\usepackage[margin=2.5cm]{geometry}
\usepackage{amsmath,amssymb,amsthm}
\usepackage{booktabs}
\usepackage{hyperref}
\emergencystretch=3em

\newcommand{\vacL}{\langle\mathrm{vacL}|}
\newcommand{\nb}[1]{\texttt{#1}}

\theoremstyle{definition}
\newtheorem{result}{Result}
\newtheorem{conjecture}{Conjecture}
\theoremstyle{remark}
\newtheorem*{remark}{Remark}

\title{CG tests for $[3,1,0]^{\otimes2}$ and $[3,2,0]^{\otimes2}$:\\
which highest-weight states do the functionals reach after untwisting?}
\author{Notes accompanying \nb{xxx-310-cg-test.wb}}
\date{2026-09-23}

\begin{document}
\maketitle

\begin{abstract}
We repeat the adjoint CG test of \nb{xxx-210-cg-test.tex} for $[3,1,0]^{\otimes2}$, $[3,2,0]^{\otimes2}$ and $[4,2,0]^{\otimes2}$, with two covector families: the one from \nb{L2-mathcalA.tex} ($\langle v_{\rm tot}|\otimes\vacL$), and the one from the determinant of \nb{L2-det-mathcalA.tex} ($\langle v_{\rm tot}|\otimes\langle v_1|$). For $[3,1,0]^{\otimes2}$ the families reach 8 of the 10 highest-weight states; $[4,2,2]$ and $[3,3,2]$ are missed. For $[3,2,0]^{\otimes2}$ they reach 9 of 10; only $[4,3,3]$ is missed. For $[4,2,0]^{\otimes2}$ they reach 15 of 19, the rule below having predicted the four missed states. The determinant family never adds a state. In all four reps tested, including the adjoint, the site-1 covector vanishes identically, at finite twist, exactly at site-1 weights with $\mu_3>\lambda_2$, and the missed irreps are exactly those with $\nu_3>\lambda_2$. Away from the missed irreps the overlap map is injective, so a highest-weight state is reconstructed uniquely from its overlaps with the generating functions. All results are exact and symbolic.
\end{abstract}

\tableofcontents

%======================================================================
\section{Setup}

\paragraph{Single site.}
One site of $V=V_\lambda$, $\lambda=[\lambda_1,\lambda_2,0]$ with $\lambda_2\le2$, in the conventions of \nb{L1-mathcalA.tex}: Gelfand--Tsetlin basis with the highest-weight vector $e_1$ first, $E_{ij}$ with $i<j$ raising, $T_{ij}(u)=(u-\theta_1)\delta_{ij}-hE_{ji}$, $t_a=\sum_{|I|=a}z_I\operatorname{qminor}_{I,I}$, and $\nu_1(u)=u-\theta_1-h\lambda_1$. The fused transfer matrices are
\begin{equation}
t_{\{1,1\}}(u)=\frac{t_2(u)}{\nu_1(u)\nu_1(u-h)},\qquad
t_{\{2,2\}}(u)=\frac{t_2(u)\,t_2(u+h)-t_3(u+h)\,t_1(u)}{\nu_1(u)\,\nu_1(u-h)\,\nu_1(u+h)\,\nu_1(u)} .
\end{equation}
Everything is exact and symbolic in $z_k,\theta_1,h,p_1,p_2,k_1,k_2$. The rep is a switch, \nb{lamSwitchC}; $\{2,1,0\}$ reproduces every count of \nb{xxx-210-cg-test.wb}.

\paragraph{The two covector families.}
By \nb{L2-mathcalA.tex} (Result~4) and \nb{L2-det-mathcalA.tex} (Result~2), with $d=\lambda_1-\lambda_2$,
\begin{equation}
\langle v_{\rm tot}|=\langle v_1|\sum_{m=0}^{\lambda_2}(-1)^m\binom{\lambda_2}{m}k_2^m\,t_{\{m,m\}}(\theta_1),\qquad
\langle v_1|=(1-k_1z_1)^d\,e_1^Tg_1,\quad g_1=e^{p_1'E_{12}}e^{p_2E_{13}}e^{p_3E_{23}},
\end{equation}
with $p_1'=p_1(1-k_1z_2)/(1-k_1z_1)$, $p_3=p_2(z_3-z_1)/(p_1(z_2-z_3))$. The families are
\begin{align}
\text{old:}\quad &\langle v_{\rm tot}|\otimes\vacL, && \vacL=e_1^Tg_0,\ g_0=e^{p_1E_{12}}e^{p_2E_{13}}e^{p_3E_{23}},\\
\text{determinant:}\quad &\langle v_{\rm tot}|\otimes\langle v_1| .
\end{align}
For $\lambda_2=1$ the sum is $1-k_2t_{\{1,1\}}$; for $\lambda_2=2$ it is $1-2k_2t_{\{1,1\}}+k_2^2t_{\{2,2\}}$.

\paragraph{Undoing the site-2 rotation.}
Both $g_0$ and $g_1$ lie in $N_+$, the group generated by the raising $E_{ij}$, $i<j$. Undoing them globally gives
\begin{align}
(\langle v_{\rm tot}|\otimes\vacL)(g_0^{-1}\otimes g_0^{-1})&=\langle v_{\rm new}|\otimes e_1^T, & v_{\rm new}&=v_{\rm tot}\,g_0^{-1},\\
(\langle v_{\rm tot}|\otimes\langle v_1|)(g_1^{-1}\otimes g_1^{-1})&=\langle v'_{\rm new}|\otimes e_1^T, & v'_{\rm new}&=(1-k_1z_1)^d\,v_{\rm tot}\,g_1^{-1}.
\end{align}
Since $v_1g_1^{-1}=(1-k_1z_1)^de_1^T$ exactly, $v'_{\rm new}=(1-k_1z_1)^{2d}\,e_1^T\sum_m(-1)^m\binom{\lambda_2}{m}k_2^m\,g_1t_{\{m,m\}}g_1^{-1}$: the $E_{12}$ shift of the old family is absent.

\begin{result}[exact]\label{res:nplus}
$g\otimes g$ fixes every highest-weight vector of $V\otimes V$ for $g\in N_+$ (checked for $g_1$ on all of them). Hence the global rotation does not change any pairing with a highest-weight state: $\langle c\otimes e_1^Tg|\,w\rangle=\langle cg^{-1}\otimes e_1^T|\,w\rangle$ for every highest-weight $w$.
\end{result}

\paragraph{Pairing and counting.}
A covector $c\otimes e_1^T$ pairs with a highest-weight vector $w$ of total weight $\nu$ only through the site-1 component of weight $\mu=\nu-\lambda$. Highest-weight vectors are computed as the common kernel of the total $E_{12}$, $E_{23}$ in each weight space of $V\otimes V$. We untwist, $z_k=e^{\varphi_k\varepsilon}$, and take the leading order in $\varepsilon$ separately in each weight space of $V$ (series to $\varepsilon^{10}$; a flag checks that no nonzero component is lost by truncation, and it is false in every run). Ranks are exact: clear a common denominator and collect the coefficient vectors of all monomials in $(k_1,k_2,p_1,p_2,\varphi_k,\theta_1,h)$. They are cross-checked by 25 random rational evaluations.

%======================================================================
\section{Per-weight leading orders}

Leading order $\varepsilon^n$ and rank of each family in each weight space of $V$; ``--'' means the component vanishes identically, at finite twist.

\begin{result}[{exact, $[3,1,0]$, $\dim V=15$}]\label{res:w310}
\begin{center}
\begin{tabular}{lccccccccccccc}\toprule
$\mu$ & 310 & 301 & 220 & 211 & 202 & 130 & 121 & 112 & 103 & 031 & 022 & 013 & total\\
$\dim V_\mu$ & 1 & 1 & 1 & 2 & 1 & 1 & 2 & 2 & 1 & 1 & 1 & 1 & 15\\\midrule
$v_{\rm new}$: order & 0 & 1 & 1 & 2 & -- & 2 & 3 & -- & -- & 4 & -- & -- & \\
\phantom{$v_{\rm new}$:} rank & 1 & 1 & 1 & 2 & 0 & 1 & 2 & 0 & 0 & 1 & 0 & 0 & 9\\\midrule
$v'_{\rm new}$: order & 0 & 1 & -- & 2 & -- & -- & -- & -- & -- & -- & -- & -- & \\
\phantom{$v'_{\rm new}$:} rank & 1 & 1 & 0 & 1 & 0 & 0 & 0 & 0 & 0 & 0 & 0 & 0 & 3\\\bottomrule
\end{tabular}
\end{center}
\end{result}

\begin{result}[{exact, $[3,2,0]$, $\dim V=15$}]\label{res:w320}
\begin{center}
\begin{tabular}{lccccccccccccc}\toprule
$\mu$ & 320 & 311 & 302 & 230 & 221 & 212 & 203 & 131 & 122 & 113 & 032 & 023 & total\\
$\dim V_\mu$ & 1 & 1 & 1 & 1 & 2 & 2 & 1 & 1 & 2 & 1 & 1 & 1 & 15\\\midrule
$v_{\rm new}$: order & 0 & 1 & 2 & 1 & 2 & 3 & -- & 3 & 4 & -- & 5 & -- & \\
\phantom{$v_{\rm new}$:} rank & 1 & 1 & 1 & 1 & 2 & 2 & 0 & 1 & 2 & 0 & 1 & 0 & 12\\\midrule
$v'_{\rm new}$: order & 0 & 1 & 2 & -- & 2 & 3 & -- & -- & 4 & -- & -- & -- & \\
\phantom{$v'_{\rm new}$:} rank & 1 & 1 & 1 & 0 & 1 & 1 & 0 & 0 & 1 & 0 & 0 & 0 & 6\\\bottomrule
\end{tabular}
\end{center}
\end{result}

For comparison, the adjoint (\nb{xxx-210-cg-test.tex}, Result~1) has $v_{\rm new}$ of rank 6 out of 8, identically zero at $(1,0,2)$ and $(0,1,2)$.

\paragraph{Observations.}
\begin{itemize}
\item In every rep, where $v_{\rm new}$ is nonzero its rank equals $\dim V_\mu$. The rank is lost only through exact zeros.
\item \textbf{$v_{\rm new}$ vanishes identically exactly at the weights with $\mu_3>\lambda_2$}: $\mu_3=2$ for the adjoint, $\mu_3\in\{2,3\}$ for $[3,1,0]$, $\mu_3=3$ for $[3,2,0]$. Each $t_{\{m,m\}}$ insertion allows one more unit of $\mu_3$.
\item $v'_{\rm new}$ is supported only on weights $\lambda-a\,\alpha_{23}-b\,\alpha_{13}$ with $a,b\ge0$, never on those needing a net $\alpha_{12}$ step (in all three reps). In multiplicity-2 spaces its rank is 1. This is consistent with $v'_{\rm new}$ lacking the $E_{12}$ shift, but it has not been derived.
\end{itemize}

%======================================================================
\section{Highest-weight states reached}

\begin{result}[exact]\label{res:hw}
Rank of the pairing of each family (per-weight leading order) with the highest-weight vectors of $V\otimes V$:
\begin{center}
\begin{tabular}{lcccc@{\qquad}lcccc}\toprule
\multicolumn{5}{c}{$[3,1,0]^{\otimes2}$} & \multicolumn{5}{c}{$[3,2,0]^{\otimes2}$}\\
irrep & \# & old & det & comb. & irrep & \# & old & det & comb.\\\midrule
$[6,2,0]$ & 1 & 1 & 1 & 1 & $[6,4,0]$ & 1 & 1 & 1 & 1\\
$[6,1,1]$ & 1 & 1 & 1 & 1 & $[6,3,1]$ & 1 & 1 & 1 & 1\\
$[5,3,0]$ & 1 & 1 & 0 & 1 & $[6,2,2]$ & 1 & 1 & 1 & 1\\
$[5,2,1]$ & 2 & 2 & 1 & 2 & $[5,5,0]$ & 1 & 1 & 0 & 1\\
$[4,4,0]$ & 1 & 1 & 0 & 1 & $[5,4,1]$ & 2 & 2 & 1 & 2\\
$[4,3,1]$ & 2 & 2 & 0 & 2 & $[5,3,2]$ & 2 & 2 & 1 & 2\\
$[4,2,2]$ & 1 & \textbf{0} & 0 & \textbf{0} & $[4,4,2]$ & 1 & 1 & 1 & 1\\
$[3,3,2]$ & 1 & \textbf{0} & 0 & \textbf{0} & $[4,3,3]$ & 1 & \textbf{0} & 0 & \textbf{0}\\\midrule
total & 10 & 8 & 3 & 8 & total & 10 & 9 & 6 & 9\\\bottomrule
\end{tabular}
\end{center}
The missed irreps need site-1 weights $(1,1,2)$, $(0,2,2)$ for $[3,1,0]$ and $(1,1,3)$ for $[3,2,0]$, exactly where $v_{\rm new}$ vanishes identically. The determinant family adds nothing to the old one.
\end{result}

\paragraph{Summary over all three reps.}
\begin{center}
\begin{tabular}{lcccl}\toprule
rep & \# hw states & reached (old = combined) & det alone & missed\\\midrule
$[2,1,0]$ & 6 & 5 & 3 & $[2,2,2]$\\
$[3,1,0]$ & 10 & 8 & 3 & $[4,2,2]$, $[3,3,2]$\\
$[3,2,0]$ & 10 & 9 & 6 & $[4,3,3]$\\
$[4,2,0]$ & 19 & 15 & 6 & $[6,3,3]$, $[5,4,3]^{\times2}$, $[4,4,4]$\\\bottomrule
\end{tabular}
\end{center}

\begin{conjecture}[{four reps; $[4,2,0]$ was a prediction}]\label{conj:rule}
With $\lambda_3=0$, $v_{\rm new}\equiv0$ exactly at the site-1 weights with $\mu_3>\lambda_2$. The families reach every highest-weight state of $V_\lambda\otimes V_\lambda$ except those in irreps with $\nu_3>\lambda_2$.
\end{conjecture}
For the adjoint this says that the singlet is missed; for $[3,1,0]$ it is not only the singlet-like states ($[3,1,0]^{\otimes2}$ has no singlet), and for $[3,2,0]$ the second insertion $t_{\{2,2\}}$ recovers all of $\nu_3=2$.

\begin{remark}
The zeros at $\mu_3>\lambda_2$ are now derived (\nb{xxx-210-cg-test.tex}, section ``Why the rule $\nu_3\le\lambda_2$''). Every family here has the form $e_1^T\,n\,X\,n^{-1}$ or $e_1^Te^{aE_{12}}X$, with $n\in N_+$ and $X=\sum_{m\le\lambda_2}c_m t_{\{m,m\}}$. By covariance $nXn^{-1}$ is the same operator with the rotated twist $nGn^{-1}$. $e_1^T$ and $e_1^Te^{aE_{12}}$ have $\mu_3=0$, and each quantum minor, being weight-homogeneous with $|I|=|J|$, moves at most one unit into slot~3 \emph{for any twist}. So $\mu_3\le\lambda_2$. This is checked exactly for the adjoint with a completely general $3\times3$ twist. The vanishing $(1,3)$ entry of $G_1$ is therefore not the cause for these families.
\end{remark}

%======================================================================
\section{$[4,2,0]^{\otimes2}$: a prediction of the rule}\label{sec:420}

Run with \nb{lamSwitchC}$=\{4,2,0\}$ ($\dim V=27$, $\dim V\otimes V=729$; the first rep with $d=2$ and $\lambda_2=2$ together), all three families, series to $\varepsilon^{10}$ (truncation flags false). The covector forms were re-verified numerically for this rep in \nb{L2-det-mathcalA.tex}.

\begin{result}[exact]\label{res:420}
$v_{\rm new}$ vanishes identically exactly at the site-1 weights with $\mu_3\in\{3,4\}$, i.e.\ $\mu_3>\lambda_2$, and is nonzero at every weight with $\mu_3\le2$. Highest-weight states reached:
\begin{center}
\begin{tabular}{lccccc}\toprule
irrep & \# hw & old & det & $w$-vacuum & all three\\\midrule
$[8,4,0]$ & 1 & 1 & 1 & 1 & 1\\
$[8,3,1]$ & 1 & 1 & 1 & 1 & 1\\
$[8,2,2]$ & 1 & 1 & 1 & 1 & 1\\
$[7,5,0]$ & 1 & 1 & 0 & 0 & 1\\
$[7,4,1]$ & 2 & 2 & 1 & 1 & 2\\
$[7,3,2]$ & 2 & 2 & 1 & 1 & 2\\
$[6,6,0]$ & 1 & 1 & 0 & 0 & 1\\
$[6,5,1]$ & 2 & 2 & 0 & 0 & 2\\
$[6,4,2]$ & 3 & 3 & 1 & 1 & 3\\
$[6,3,3]$ & 1 & \textbf{0} & 0 & 0 & \textbf{0}\\
$[5,5,2]$ & 1 & 1 & 0 & 0 & 1\\
$[5,4,3]$ & 2 & \textbf{0} & 0 & 0 & \textbf{0}\\
$[4,4,4]$ & 1 & \textbf{0} & 0 & 0 & \textbf{0}\\\midrule
total & 19 & 15 & 6 & 6 & 15\\\bottomrule
\end{tabular}
\end{center}
The four missed states are exactly those with $\nu_3>\lambda_2=2$, as Conjecture~\ref{conj:rule} predicts. Where the old family is nonzero it has full rank, including the multiplicity-3 space $[6,4,2]$. At finite twist the overlap map has rank 15 on all 19 highest-weight vectors (old family alone, and all three together); the $w$-vacuum family alone has rank 6.
\end{result}

%======================================================================
\section{Reconstructing a highest-weight state}

\paragraph{Question.}
Given a highest-weight state $|A\rangle$ of $V\otimes V$, do the overlaps $\langle v_{\rm new}(k_1,k_2,p,z)\otimes e_1^T|A\rangle$, known as functions of all parameters, determine $|A\rangle$?

\paragraph{Answer.}
The map $|A\rangle\mapsto\langle v_{\rm new}(\cdot)\otimes e_1^T|A\rangle$ is linear, so $|A\rangle$ is determined (normalization included) exactly when the map is injective on the relevant space. For a state of definite weight $\nu$ that space is the multiplicity space of highest-weight vectors of weight $\nu$, and Result~\ref{res:hw} shows full rank there for every $\nu$ that is not missed, including the multiplicity-2 cases. For states mixing several weights we check the whole space at once.

\begin{result}[exact, finite twist]\label{res:recon}
Rank of the overlap map on the span of \emph{all} highest-weight vectors of $V\otimes V$:
\begin{center}
\begin{tabular}{lccc}\toprule
rep & \# hw vectors & rank, old family & rank, old + determinant\\\midrule
$[2,1,0]$ & 6 & 5 & 5\\
$[3,1,0]$ & 10 & 8 & 8\\
$[3,2,0]$ & 10 & 9 & 9\\\bottomrule
\end{tabular}
\end{center}
The kernel is exactly the span of the missed irreps. So a highest-weight state with no component in the irreps with $\nu_3>\lambda_2$ is fixed uniquely by its overlaps; a component in those irreps is invisible.
\end{result}

\begin{remark}
In the untwisted limit the statement holds weight by weight (Result~\ref{res:hw}, from the per-weight leading orders); the mixed-weight check is at finite twist only. For the adjoint, separating the two octets needs the $k_1$-dependence: with $k_1$ frozen the rank of that sector drops to 1 (\nb{xxx-210-cg-test.tex}, Result~3). Non-highest-weight states are not reconstructed directly (with site~2 fixed at $e_1^T$ the covectors span at most $\dim V$ directions), but they follow from the highest-weight states by lowering, which is all the CG problem needs.
\end{remark}

%======================================================================
\section{Status and open points}
\begin{itemize}
\item \textbf{Exact and symbolic:} Results~\ref{res:nplus}--\ref{res:recon}.
\item \textbf{Inherited, not re-checked here:} the covector forms $\langle v_{\rm tot}|\otimes\vacL$ and $\langle v_{\rm tot}|\otimes\langle v_1|$, verified numerically in \nb{L2-mathcalA.tex} and \nb{L2-det-mathcalA.tex}.
\item \textbf{Conjecture~\ref{conj:rule}:} the half ``irreps with $\nu_3>\lambda_2$ are missed'' is now derived (remark after it). The half ``everything else is reached'' is still observed only, on four reps; $[4,2,0]$ was a prediction.
\item \textbf{Open:}
  \begin{itemize}
  \item a derivation of the zeros at $\mu_3>\lambda_2$: done (remark after Conjecture~\ref{conj:rule}; details in \nb{xxx-210-cg-test.tex});
  \item functionals that reach $\nu_3>\lambda_2$;
  \item $\lambda_2=3$ ($t_{\{3,3\}}$), $\lambda_3\neq0$, and different reps on the two sites.
  \end{itemize}
\end{itemize}

\paragraph{Pitfalls met.}
\begin{itemize}
\item The adjoint cells 6, 7, 10 of \nb{xxx-210-cg-test.wb} (octet parameter freezing, mechanism, singlet checks) are adjoint-specific and were dropped here.
\item The saved outputs of the notebook are the last run, $[3,2,0]$. The $[2,1,0]$ and $[3,1,0]$ numbers above were produced by the same cells with \nb{lamSwitchC} changed.
\item The notebook shares its kernel with the other XXX-CG notebooks, and \nb{lamSwitchC} persists across runs.
\end{itemize}

%======================================================================
\section*{Notebook map (\nb{xxx-310-cg-test.wb})}
\begin{center}
\begin{tabular}{rl}\toprule
cell & content\\\midrule
1 & summary and three-rep table\\
2 & GT machinery\\
3 & switch \nb{lamSwitchC}; exact $t_{\{m,m\}}$, $\langle v_{\rm tot}|$, $g_0$, $v_{\rm new}$\\
4 & untwisting, per-weight leading orders and ranks of $v_{\rm new}$; truncation flag\\
5 & $V\otimes V$, highest-weight vectors, pairing ranks of the old family\\
6--7 & determinant family $v'_{\rm new}$; exact zeros; old / det / combined counts (Results~\ref{res:w310}--\ref{res:hw})\\
8 & rank of the overlap map on all highest-weight vectors (Result~\ref{res:recon})\\\bottomrule
\end{tabular}
\end{center}
To rerun for another rep, set \nb{lamSwitchC} and run cells 3--8 (seconds to a few minutes).

\end{document}
